\chapter{腔模量子化与真空涨落}
\label{chap:cavity qed}

\chapref{chap:waveguide}中，壁面的边界条件决定了电磁场可以怎样振荡。现在让腔中一个原子从激发态跃迁到基态，并把能量交给其中一个模式。原子在什么位置、跃迁偶极朝哪个方向，决定了这一过程有多强。先由经典场能确定一个光子所对应的场强，便能把这些几何因素带进寿命和频谱的计算。\footnote{场的正则量子化可对照D. Tong的\href{https://www.damtp.cam.ac.uk/user/tong/qft/qfthtml/S6.html}{Quantum Field Theory，第6章}。单模腔与两能级原子的耦合见D. A. Steck的\href{https://atomoptics-nas.uoregon.edu/~dsteck/teaching/quantum-optics/}{Quantum and Atom Optics}。}

\section{从圆柱腔模到单个光子}
\label{sec:single cavity mode}

仍取半径$a$、长$d$的真空圆柱腔，先将壁视为理想导体。TM$_{010}$模的电场沿轴向，取实模函数$\bm u$并将其峰值归一为1，
\[
    \bm u(\bm r)=\uvec z J_0(x_{01}\rho/a),\qquad
    J_0(x_{01})=0,\quad x_{01}\simeq2.405,\qquad \omega_c=cx_{01}/a .
\]
定义模体积$V_{\mathrm m}=\int_V\abs{\bm u}^2\d V$。由$J_0$在侧壁上为零，径向积分给出
\begin{equation}
    V_{\mathrm m}=2\pi d\int_0^a \rho J_0^2(x_{01}\rho/a)\d\rho
       =\pi a^2d J_1^2(x_{01})\simeq0.270\,\pi a^2d .
    \label{eqn:tm010 mode volume}
\end{equation}
因此用轴上峰值估算全腔电场能时，应乘这个模体积。离轴电场的减弱已经包含在积分中。

\paragraph{场能给出正则坐标}
在无源腔内取Coulomb规范和$\Phi=0$，以实函数$q(t)$为这个模式的广义坐标，令
\[
    \bm A(\bm r,t)=\frac{q(t)\bm u(\bm r)}{\sqrt{\varepsilon_0 V_{\mathrm m}}},
    \qquad \bm E=-\pv{\bm A}t,\qquad \bm B=\curl\bm A .
\]
模函数满足$\curl(\curl\bm u)=(\omega_c/c)^2\bm u$，且$\uvec n\times\bm u=0$。分部积分消去壁面项，得
$\int_V\abs{\curl\bm u}^2\d V=(\omega_c/c)^2 V_{\mathrm m}$。对瞬时实场积分，电场能和磁场能分别为
\[
    U_{\electric}=\frac{\varepsilon_0}{2}\int_V\bm E^2\d V=\frac{\dot q^2}{2},
    \qquad
    U_{\magnetic}=\frac{1}{2\mu_0}\int_V\bm B^2\d V=\frac{\omega_c^2q^2}{2}.
\]
故$L=U_\electric-U_\magnetic$给出正则动量$p=\dot q$，Hamilton量为$H=(p^2+\omega_c^2q^2)/2$。若经典轴上场的峰值为$E_0$，总储能等于电场能最大值$\varepsilon_0 V_{\mathrm m}\abs{E_0}^2/2$，也与\chapref{chap:waveguide}的复振幅计算相符。

将$q,p$提升为算符，使$[\hat q,\hat p]=\i\hbar$。写成满足$[\hat a,\hat a^\dagger]=1$的产生、湮灭算符，
\[
    \hat q=\sqrt{\frac{\hbar}{2\omega_c}}(\hat a+\hat a^\dagger),\qquad
    \hat p=-\i\sqrt{\frac{\hbar\omega_c}{2}}(\hat a-\hat a^\dagger),\qquad
    \hat H=\hbar\omega_c\bigp{\hat a^\dagger\hat a+\tfrac12}.
\]
由于电场正比于$-\hat p$，在选定的模式相位下有
\begin{equation}
    \mel 0{\hat{\bm E}(\bm r)}1
      =\i\sqrt{\frac{\hbar\omega_c}{2\varepsilon_0 V_{\mathrm m}}}\,
       J_0(x_{01}\rho/a)\uvec z .
    \label{eqn:tm010 vacuum matrix element}
\end{equation}
这里$\ket1=\hat a^\dagger\ket0$。真空态的平均电场为零，轴上这个模式的电场均方根却等于式\eqref{eqn:tm010 vacuum matrix element}的模。原子沿轴移动时矩阵元不变，沿半径靠近侧壁时则趋于零。

\section{边界附近的发射率}
\label{sec:cavity spontaneous emission}

设静止的两能级原子有向下跃迁角频率$\omega_a>0$，跃迁偶极矩为$\bm d_{ge}=\mel g{\hat{\bm d}}e$，原子位于$\bm r_a$。自由空间中，电偶极近似给出单个通道的发射率
\[
    \Gamma_0=\frac{\omega_a^3\abs{\bm d_{ge}}^2}{3\pi\varepsilon_0\hbar c^3}.
\]
《量子力学》卷第8章由两个横向偏振的求和推导了这个结果。改变边界时，电偶极相互作用$\hat H_{\rm int}=-\hat{\bm d}\cdot\hat{\bm E}(\bm r_a)$依然成立，变化的是每个模式的场和允许发射的末态。

取真空区域无损电磁场的正频率本征模$\bm f_\nu$，归一化为$\int \bm f_\nu\conj\cdot\bm f_{\nu'}\d V=\delta_{\nu\nu'}$。在大盒极限中用Fermi黄金规则对单光子末态求和，有
\[
    \Gamma=\frac{\pi\omega_a}{\varepsilon_0\hbar}
       \sum_\nu\abs{\bm f_\nu\conj(\bm r_a)\cdot\bm d_{ge}}^2
                \delta(\omega_a-\omega_\nu).
\]
这正是电磁Green张量的谱展开。对\secref{sec:dyadic-green}的推迟Green张量，在正频率取
$\Im[(\omega_\nu^2-(\omega+\i0)^2)^{-1}]=\pi\delta(\omega-\omega_\nu)/(2\omega)$，
可将求和写为
\begin{equation}
    \Gamma=\frac{2\omega_a^2}{\hbar\varepsilon_0 c^2}\,
    \bm d_{ge}\conj\cdot\Im\bm G(\bm r_a,\bm r_a;\omega_a)\cdot\bm d_{ge}.
    \label{eqn:decay from Green tensor}
\end{equation}
这里假设场起初处于真空，耦合较弱，且环境记忆时间短于原子布居的衰减时间。自由空间的同点虚部为$(\omega_a/6\pi c)\bm I$，代入即恢复$\Gamma_0$。

\begin{example}{镜前偶极的寿命}{mirror changes atomic lifetime}
    将原子移至距理想导体平面$h$处，保持跃迁偶极矩和跃迁频率固定。按给定取向，式\eqref{eqn:decay from Green tensor}的比值$\Gamma/\Gamma_0$恰等于同取向经典偶极的功率比$P/P_0$。在$h=\lambda_a/4$处，$\lambda_a=2\pi c/\omega_a$，往返传播相位为$\pi$。由式\eqref{eqn:mirror-dipole-rates}，
    \[
        \frac{\Gamma_\parallel}{\Gamma_0}=1+\frac{3}{2\pi^2},
        \qquad \frac{\Gamma_\perp}{\Gamma_0}=1+\frac{3}{\pi^2}.
    \]
    若这条跃迁是唯一衰减通道，寿命就分别为自由空间寿命的$0.868$倍和$0.767$倍。在$h\ll\lambda_a$的理想镜面模型中，平行偶极及其反向像偶极在远区相消，垂直偶极则与同向像偶极叠加，两个发射率分别趋于零和$2\Gamma_0$。
\end{example}

\section{腔中的寿命与真空Rabi劈裂}
\label{sec:purcell and vacuum rabi}

将同一原子置于圆柱腔内。单光子发射的耦合角频率$g$由式\eqref{eqn:tm010 vacuum matrix element}给出，选取原子和光子态的相位使$g\ge0$。令$\eta$表示位置和取向带来的因子，
\begin{equation}
    g^2=\frac{\omega_c\abs{\bm d_{ge}}^2}{2\hbar\varepsilon_0 V_{\mathrm m}}\eta,
    \qquad \eta=J_0^2(x_{01}\rho_a/a)
        \frac{\abs{\uvec z\cdot\bm d_{ge}}^2}{\abs{\bm d_{ge}}^2}.
    \label{eqn:local cavity coupling}
\end{equation}
轴上沿$z$的偶极取得最大耦合$g_{\max}$，与$z$垂直的偶极与此模耦合为零。下文假设此模孤立、近共振，且$g,\abs{\omega_a-\omega_c}$及各线宽都远小于$\omega_c$，采用旋转波近似。

\paragraph{弱耦合下的Purcell效应}
沿用式\eqref{eqn:cavity wall Q}的约定，腔模能量衰减率为$\kappa=\omega_c/Q$，这里$Q$计入壁损及输出耦合等全部损耗。腔场振幅的衰减率为$\kappa/2$。原子通过其余通道的布居衰减率记为$\Gamma_{\rm bg}$。忽略纯退相干，在以$\omega_a$旋转的参考系中，原子激发态振幅$c_e$与腔内单光子振幅$c_1$满足
\begin{align}
    \dot c_e&=-\tfrac12\Gamma_{\rm bg}c_e-\i g c_1, &
    \dot c_1&=(\i\Delta-\tfrac12\kappa)c_1-\i g c_e,
    \qquad \Delta=\omega_a-\omega_c .
    \label{eqn:atom cavity amplitudes}
\end{align}
这些振幅描述能量尚留在原子和所选腔模内的单激发分量。若$\kappa\gg g,\Gamma_{\rm bg}$，腔场在原子布居明显改变前就已衰减。在最初的腔瞬态之后，取$\dot c_1\simeq0$，代回原子方程，由$\abs{c_e}^2$的衰减得到
\[
    \Gamma_{\rm tot}=\Gamma_{\rm bg}+\Gamma_c,\qquad
    \Gamma_c=\frac{g^2\kappa}{\Delta^2+(\kappa/2)^2}.
\]
因而单激发原子的寿命为$1/\Gamma_{\rm tot}$。只有当其他通道与自由空间相同，才可置$\Gamma_{\rm bg}=\Gamma_0$。将腔通道单独与自由空间的同一跃迁比较，得到
\begin{equation}
    \frac{\Gamma_c}{\Gamma_0}
       =\frac{3}{4\pi^2}\frac{\lambda_a^3 Q}{V_{\mathrm m}}\,
        \frac{\eta}{1+(2\Delta/\kappa)^2}.
    \label{eqn:purcell rate}
\end{equation}
共振时腔通道的发射率为$4g^2/\kappa$。高$Q$把可用光子态集中在较窄的频带，小模体积增大归一化场强。这两个效应在式\eqref{eqn:purcell rate}中以$Q/V_{\mathrm m}$出现。

\paragraph{共振下的两个本征频率}
减小腔损耗，腔光子便有时间重新激发原子。以$\ket{e,0},\ket{g,1}$为基底、以$\ket{g,0}$的能量为零点，封闭系统的Hamilton量为
\[
    \frac{\hat H_1}{\hbar}=
    \begin{pmatrix}\omega_a&g\\g&\omega_c\end{pmatrix}.
\]
当$\omega_a=\omega_c$时，对称和反对称叠加态$(\ket{e,0}\pm\ket{g,1})/\sqrt 2$的角频率为$\omega_c\pm g$。从$\ket{e,0}$出发，原子布居按$\cos^2(gt)$变化。频域中两个允许的跃迁相距$2g$，称为真空Rabi劈裂。

考虑损耗后，在矩阵的两个对角元分别加上$-\i\Gamma_{\rm bg}/2$和$-\i\kappa/2$。共振时复频率为
\begin{equation}
    \tilde\omega_\pm=\omega_c-\frac{\i}{4}(\kappa+\Gamma_{\rm bg})
       \pm\sqrt{g^2-\frac{(\kappa-\Gamma_{\rm bg})^2}{16}} .
    \label{eqn:damped vacuum rabi}
\end{equation}
当$g>\abs{\kappa-\Gamma_{\rm bg}}/4$时，两极点已有不同实部。要在弱探测频谱上清楚看见双峰，还需两支模式都耦合到探测通道，峰间距大于各支的线宽和仪器分辨宽度。在$g\gg\kappa,\Gamma_{\rm bg}$的极限，每支峰的功率半高全宽约为$(\kappa+\Gamma_{\rm bg})/2$。

\begin{example}{改变同一模式的输出损耗}{cavity lifetime and splitting}
    设原子位于轴上、偶极沿$z$，固定$\Gamma_{\rm bg}=g/100$。令原子与腔共振，调节输出耦合，在小扰动下保持腔模形状和$g$固定。若$\kappa=40g$，腔带来的发射率为$0.10g$，布居寿命约为$1/(0.11g)$。将失谐调至$\abs\Delta=\kappa/2$，此通道减半，寿命增至约$1/(0.06g)$。
    若调回共振并使$\kappa=0.20g$，式\eqref{eqn:damped vacuum rabi}给出的极点间距约为$2.00g$，每支线宽约为$0.105g$。此时原子与光子相干交换能量，低损耗使两个联合模式可以在频谱上分辨。
\end{example}

在同样的初态$c_e(0)=1,\ c_1(0)=0$下，可以直接计算这两种测量。图\ref{fig:cavity observables}左图由式\eqref{eqn:atom cavity amplitudes}求出激发态布居。右图取低损耗共振腔，测从腔输出口离开的光子。记$\delta=\omega-\omega_c$，输出谱正比于$\abs{\int_0^\infty c_1(t)\e^{\i\delta t}\d t}^2$。将两个振幅方程作单边Fourier变换，得到归一化前的谱形
\[
    S_{\rm out}(\delta)\ \propto\
    \frac{g^2}{\bigabs{(\kappa/2-\i\delta)(\Gamma_{\rm bg}/2-\i\delta)+g^2}^2}.
\]
\begin{figure}[htbp]
    \centering
    \includegraphics[width=.92\linewidth]{figures/cavity_decay_spectrum.pdf}
    \caption{同一跃迁的激发态布居和腔输出谱。左图取$\kappa=40g$，右图取$\kappa=0.20g$且腔与原子共振，两图均取$\Gamma_{\rm bg}=0.01g$。}
    \label{fig:cavity observables}
\end{figure}

\section{平行导体板的Casimir力}
\label{sec:parallel plates Casimir}

单个腔模在基态也有$\hbar\omega/2$的能量。改变板距会改变整组本征频率，可以通过板间压力测得基态能量差对几何的依赖。\footnote{H. B. G. Casimir，\href{https://web.mit.edu/kardar/www/research/seminars/Casimir/Casimir1948.pdf}{On the Attraction Between Two Perfectly Conducting Plates} (1948)。}考虑零温下两块面积为$S$、相距$a$的平行理想导体，板间为真空，忽略边缘。平行于板面的$\bm k_\parallel$连续，法向波数取$n\pi/a$，故
\[
    \omega_n=c\sqrt{k_\parallel^2+(n\pi/a)^2},\qquad n=0,1,2,\ldots .
\]
对$n\ge1$有两种偏振。对$n=0$，沿板法向的电场允许存在，另一偏振的切向电场在两块板上均须为零，故只剩一种偏振。

用同一长度参数$\ell>0$为所有零点能乘上截止因子$\exp(-\ell\omega/c)$。保留两板相互作用引起的能量差时，须减去板距无限大、相同体积内的真空能。这样内外相同的体能密度不给板施加净力。设$b=\pi/a$，对横向波矢以$k=\omega/c$换元，$k_\parallel\d k_\parallel=k\d k$，得
\begin{align*}
    \frac{E_{\rm int}(a,\ell)}{S}
       &=\frac{\hbar c}{2\pi}\biggl[
           \sum_{n=0}^\infty{}'F_\ell(nb)
           -\frac{1}{b}\int_0^\infty F_\ell(t)\d t\biggr],\\
    F_\ell(t)&=\int_t^\infty k^2\e^{-\ell k}\d k .
\end{align*}
求和符号的撇表示$n=0$项取半权，已将两种偏振的公共因子提到求和之外。连续极限的法向模密度为$1/b$。在截止因子相同的条件下作差，可用Euler--Maclaurin公式求其有限部分，
\[
    \sum_{n=0}^\infty{}'F_\ell(nb)-\frac{1}{b}\int_0^\infty F_\ell(t)\d t
       =-\frac{b}{12}F_\ell'(0)+\frac{b^3}{720}F_\ell'''(0)+\bigo(b^5\ell^2).
\]
由于$F_\ell'(t)=-t^2\e^{-\ell t}$，前两个所需导数为$F_\ell'(0)=0$、$F_\ell'''(0)=-2$。固定$a$取$\ell\to0^+$，得到单位面积的相互作用能及法向压力
\begin{equation}
    \frac{E_{\rm int}(a)}{S}=-\frac{\pi^2\hbar c}{720a^3},\qquad
    P=-\dd a\biggl(\frac{E_{\rm int}}S\biggr)=-\frac{\pi^2\hbar c}{240a^4}.
    \label{eqn:Casimir pressure}
\end{equation}
负号表示两块板相吸。若$a=\qty{1}{\micro\metre}$，理想模型给出$\abs P\simeq\qty{1.30e-3}{\pascal}$。将板距减半，吸引压强增大到16倍。有限导电率或有限温度下须相应改变反射边界或场模的占据数。
